\documentclass[11pt]{article}
\usepackage[T1]{fontenc}
\usepackage{lmodern}
\usepackage{amsmath,amssymb,amsthm,mathtools}
\usepackage[margin=1in]{geometry}
\usepackage{microtype}
\usepackage{booktabs,array,longtable}
\usepackage{enumitem}
\usepackage{xcolor}
\usepackage[hyphens]{url}
\usepackage[colorlinks=true,linkcolor=blue!45!black,citecolor=blue!45!black,urlcolor=blue!45!black]{hyperref}
\hypersetup{pdftitle={Local Fourier Sign Codes: Exact Obstructions, Robust Sidorenko Bounds, and Sharp Girth Forcing},pdfauthor={AI-assisted research manuscript prepared for ProveIt},pdfsubject={Finite Fourier sign codes and graph homomorphism inequalities},pdfkeywords={Sidorenko, Gowers U2, Fourier sign codes, forcing, finite abelian groups}}
\usepackage{fancyhdr}
\pagestyle{fancy}
\fancyhf{}
\fancyhead[L]{\small Local Fourier sign codes}
\fancyhead[R]{\small ProveIt research manuscript}
\fancyfoot[C]{\thepage}
\setlength{\headheight}{14pt}
\setlength{\emergencystretch}{2em}
\setlist{itemsep=3pt,topsep=5pt}
\numberwithin{equation}{section}
\newtheorem{theorem}{Theorem}[section]
\newtheorem{proposition}[theorem]{Proposition}
\newtheorem{lemma}[theorem]{Lemma}
\newtheorem{corollary}[theorem]{Corollary}
\theoremstyle{definition}
\newtheorem{definition}[theorem]{Definition}
\newtheorem{example}[theorem]{Example}
\theoremstyle{remark}
\newtheorem{remark}[theorem]{Remark}
\newcommand{\E}{\mathbb E}
\newcommand{\F}{\mathbb F}
\newcommand{\Z}{\mathbb Z}
\newcommand{\R}{\mathbb R}
\newcommand{\Q}{\mathbb Q}
\newcommand{\C}{\mathbb C}
\newcommand{\T}{\mathbb T}
\newcommand{\one}{\mathbf 1}
\newcommand{\supp}{\operatorname{supp}}
\newcommand{\sgn}{\operatorname{sgn}}
\newcommand{\rank}{\operatorname{rank}}
\newcommand{\Span}{\operatorname{span}}
\newcommand{\Cut}[1]{\left\|#1\right\|_{\square}}
\newcommand{\Fl}{\mathcal F}
\newcommand{\Err}{\mathcal E}
\newcommand{\Rel}{\mathcal R}
\newcommand{\Hstar}{H_{\!*}}
\newcommand{\eps}{\varepsilon}
\newcommand{\abs}[1]{\left|#1\right|}
\newcommand{\norm}[1]{\left\|#1\right\|}

\title{\textbf{Local Fourier Sign Codes}\\[5pt]
\large Exact Obstructions, Robust Sidorenko Bounds,\\
and Sharp Girth Forcing}
\author{Research manuscript prepared for the ProveIt project\\
\small AI-assisted mathematical development; not independently refereed}
\date{7 October 2026}

\begin{document}
\maketitle
\begin{abstract}
For a real even function on a finite abelian group, we classify the
Fourier-sign obstruction to proving bipartite homomorphism inequalities
term by term. Short zero-sum relations define a binary linear code. Its
first violated constraint is exactly the first one-sided degree threshold
at which a negative conserved Fourier monomial can occur; a simple theta
graph realizes the obstruction. The shortest obstruction is computable in
a two-sheeted Cayley graph, and its length is at most the order of the
generated frequency subgroup, a bound attained for odd cyclic groups.

We then replace exact sign constraints by explicit moment defects. These
give a lower bound consisting of the random-density term, the contributions
of all simple cycles, and a controlled error. For a cyclic pattern of girth
$g$, with $c_g$ shortest cycles, this yields the sharp cut-norm forcing scale
$p(\delta+\eps)^{1/g}/(4c_g^{1/g})$. Both its exponent and its asymptotic
constant are optimal in the weighted scalar Cayley class.

For the 35-vertex, 66-edge incidence pattern specified in the
\texttt{openai/math} Sidorenko manuscript, we prove a cubic-defect bound
with an exact $33p^{62}\|w-p\|_{U^2}^4$ surplus. Thus a scalar Cayley
counterexample for this pattern would require a quantitatively nonzero
mass of sign-frustrated frequency triples. We also give its complete
one-character density polynomial and reproducible exact-arithmetic
checks. The results do not assume the correctness of the repository's
counterexample construction, do not assert a counterexample of our own,
and do not settle the unrestricted Sidorenko problem. Priority of the
combined refinements requires further literature and expert review.
\end{abstract}

\vspace{3pt}
\noindent\textbf{Keywords.} Fourier analysis; graph homomorphisms; Sidorenko inequalities;
Gowers uniformity; signed zero-sum relations; binary codes; forcing.

\clearpage
\setcounter{tocdepth}{1}
\tableofcontents
\clearpage

\section{Research setting, contributions, and scope}
\label{sec:intro}

\subsection{The intersection of the two repositories}

The \texttt{openai/math} catalogue includes a manuscript titled
\emph{A counterexample to Sidorenko's conjecture}, dated 23 September 2026.
Its finite-complex section specifies an incidence graph with 13 point
vertices, 22 face vertices, and 66 edges \cite{OpenAICatalogue,OpenAIComplex}.
Its kernel section uses finite-field rank layers, several matrix
coordinates, vertex types, and activation laws \cite{OpenAIKernel}.
These are precise source statements, not independent certification here
of the manuscript's headline claim.

The Gowers--Szemeredi development in \texttt{ProveIt} organizes finite
Fourier analysis, additive relations, uniformity, and local quantitative
estimates. Its inspected definitions use an unnormalized discrete Fourier
transform on $\Z/N\Z$, together with additive-cube and Bohr-set language
\cite{ProveItDefinitions}. A companion source audit carefully distinguishes
statement catalogues from kernel-checked proofs \cite{ProveItAudit}.
That distinction is retained throughout this article.

A useful common question is therefore not simply whether a large
construction can be repeated. It is the following.

\begin{quote}
Which local additive relations among Fourier frequencies permit negative
graph-count contributions, and how much such local negativity is required
for a global deficit?
\end{quote}

We answer an exact version for termwise Fourier positivity, and a
quantitative necessary-condition version for the actual homomorphism
inequality. The bridge to Gowers's language is particularly direct for
four-cycles: the relevant spectral fourth moment is precisely the fourth
power of the normalized $U^2$ norm.

\subsection{What is proved here}

There are four principal results. First, the safe signs on a prescribed
symmetric spectrum form the orthogonal complement of a finite binary
relation code. The first failed relation is realized as a negative flow
on a simple theta graph, giving an exact degree threshold
(Theorems~\ref{thm:threshold} and \ref{thm:code}). A finite cutoff and a
shortest-witness algorithm make the criterion effective
(Theorem~\ref{thm:bfs}).

Second, a difference of ordinary moments of $w$ and its Fourier-absolute
companion bounds the entire negative-flow mass. The resulting robust
inequality retains the contribution of \emph{every} simple cycle
(Theorem~\ref{thm:inventory}). Third, this implies a sharp girth-controlled
forcing theorem, including the exact asymptotic cut-norm constant and an
explicit exclusion of forests (Theorems~\ref{thm:forcing} and
\ref{thm:sharp}). Fourth, applying these results to the source incidence
pattern gives a cubic obstruction certificate and a sharp $U^2$ surplus
(Theorem~\ref{thm:incidence}).

The basic ingredients---Fourier orthogonality, sign characters, cycle
flows, and the even-subgraph expansion---are standard mathematical
mechanisms, not claimed inventions. The proposed contributions are their
precise local-sign classification, the explicit robust error functional,
the sharp cycle-inventory forcing bounds, and the concrete incidence
application. All are proved below, but the claim that any one formulation
is absent from all earlier literature is \emph{not} established.

\subsection{Comparison with inspected literature}

Yuqi Zhao proves Sidorenko-type inequalities for even subdivisions in
finite abelian Cayley hosts using Fourier pairing
\cite{YuqiZhao}. We rederive the required positivity directly from vertex
orthogonality, so no circuit-matrix rank assertion over a general finite
abelian group is needed. The final discussion of that preprint uses an
edge-count forcing exponent. For cyclic patterns, our shortest-cycle
argument gives the exponent $g$, counts all $c_g$ shortest cycles, and
proves optimality in the weighted class. The forest exception and the
distinction between normalized operator eigenvalues and adjacency-matrix
eigenvalues are made explicit in Section~\ref{sec:forcing}.

The positive-graph problem studied by Conlon, Lee, and Versteegen
\cite{PositiveGraphs} concerns positivity across arbitrary weighted hosts.
Here the graph and a scalar abelian convolution host are both relevant;
our termwise sign condition must not be confused with a classification
of positive graphs. Likewise, a finite relation code tests a sufficient
proof mechanism for Sidorenko, not the necessity of Sidorenko itself.

\subsection{Trust boundary}

The proof chain uses only finite Fourier duality, elementary finite sums,
linear algebra over $\F_2$, basic inequalities, and the explicit face list.
It imports no unproved repository theorem. No Lean build or kernel check
of this manuscript was performed. The exact Python checks are regression
and finite-certificate checks, not a substitute for the proofs.

In particular, the actual typed matrix kernel in
\cite{OpenAIKernel} is not identified here with a single scalar kernel
$w(x-y)$ on a uniformly measured finite abelian group. The scalar-host
conclusions below therefore neither contradict that construction nor
certify its validity. The source audit and integration notes distributed
with the article specify the inspected files and the remaining formal
verification work.

\section{Normalized Fourier flows}
\label{sec:flows}

Let $G$ be a finite abelian group with uniform probability measure, and
write its character group $\Gamma=\widehat G$ additively. For
$\xi\in\Gamma$, write $\chi_\xi:G\to\T$ for the corresponding character.
The identity of $\Gamma$ is denoted by $0$. Define
\begin{equation}
 \lambda_\xi=\widehat w(\xi)
   =\E_{x\in G}w(x)\overline{\chi_\xi(x)},
 \qquad
 w(x)=\sum_{\xi\in\Gamma}\lambda_\xi\chi_\xi(x).
 \label{eq:fourier}
\end{equation}
Throughout, $w:G\to\R$ is even: $w(-x)=w(x)$. Consequently
\begin{equation}
 \lambda_\xi\in\R,\qquad
 \lambda_{-\xi}=\lambda_\xi.
 \label{eq:real-even}
\end{equation}
Write $p=\lambda_0=\E w$, and assume $p>0$. Unless explicitly stated
otherwise, the algebraic statements allow $w$ to be signed. A
\emph{weighted scalar Cayley host} means the graphon
\begin{equation}
 W(x,y)=w(x-y),\qquad 0\leq w\leq1.
 \label{eq:host}
\end{equation}
It is symmetric and has constant degree $p$. Values on the diagonal are
included; this is a finite probability-space graphon, not necessarily a
loopless $0$--$1$ graph.

Let $H=(L\sqcup R,E)$ be a finite simple bipartite graph and set $e=|E|$.
Isolated vertices and disconnected graphs are permitted. The
homomorphism density is
\begin{equation}
 t(H,W)=\E_{(x_v)\in G^{V(H)}}
          \prod_{uv\in E}w(x_u-x_v),
 \label{eq:density}
\end{equation}
where each edge is written with $u\in L$, $v\in R$. Samples at different
vertices are independent and are allowed to coincide.

\begin{definition}[Conserved frequency labeling]
A map $\phi:E\to\Gamma$ is a \emph{Fourier flow} if
\begin{equation}
 \sum_{a\ni v}\phi(a)=0\quad\text{for every }v\in V(H).
 \label{eq:unsigned-flow}
\end{equation}
The set of such maps is $\Fl_H(\Gamma)$, and its weight is
$A(\phi)=\prod_{a\in E}\lambda_{\phi(a)}$.
\end{definition}

The unsigned appearance of \eqref{eq:unsigned-flow} is deliberate:
orient every edge from $L$ to $R$. At a left vertex all incidence signs
are positive, and at a right vertex all are negative. Either condition
is equivalent to the displayed zero sum. No division by two or assumption
on the exponent of $G$ is involved.

\begin{lemma}[Flow expansion]
\label{lem:flow}
For every finite bipartite $H$,
\begin{equation}
 t(H,W)=\sum_{\phi\in\Fl_H(\Gamma)} A(\phi).
 \label{eq:flow-expansion}
\end{equation}
The identically zero flow contributes $p^e$.
\end{lemma}
\begin{proof}
Insert \eqref{eq:fourier} at every edge. A fixed labeling contributes the
coefficient product times
\[
 \prod_{u\in L}\E_{x_u}\chi_{\sum_{a\ni u}\phi(a)}(x_u)
 \prod_{v\in R}\E_{x_v}\overline{
       \chi_{\sum_{a\ni v}\phi(a)}(x_v)}.
\]
Character orthogonality makes this expression one precisely when all
vertex sums vanish and zero otherwise. All sums are finite, so their
interchange is immediate. There is no additional $|G|$ factor: the
averages in \eqref{eq:fourier} already fix the normalization.
\end{proof}

Define the active nonzero spectrum and its signs by
\begin{equation}
 S=\{\xi\ne0:\lambda_\xi\ne0\},\qquad
 \sigma(\xi)=\sgn(\lambda_\xi)\quad(\xi\in S).
 \label{eq:spectrum}
\end{equation}
Then $S=-S$ and $\sigma(-\xi)=\sigma(\xi)$. When helpful extend
$\sigma$ to $0$ by $\sigma(0)=1$.

\begin{lemma}[Realizability of prescribed signs]
\label{lem:realizable}
Every finite symmetric $S\subseteq\Gamma\setminus\{0\}$ and every even
sign function $\sigma:S\to\{\pm1\}$ occur as the active spectrum and
signs of a weighted scalar Cayley host with any prescribed
$p\in(0,1)$.
\end{lemma}
\begin{proof}
Choose positive even numbers $a_\xi=a_{-\xi}$ with
$\sum_{\xi\in S}a_\xi\leq\min(p,1-p)$ and put
$w(x)=p+\sum_{\xi\in S}\sigma(\xi)a_\xi\chi_\xi(x)$.
Pairing inverse frequencies shows that $w$ is real and even. The triangle
inequality gives $0\leq w\leq1$. Fourier inversion gives exactly the
specified coefficients and no others.
\end{proof}

\section{Exact local sign obstructions}
\label{sec:signs}

\begin{definition}[Frustration length]
The \emph{frustration length} of $(S,\sigma)$ is
\begin{equation}
 r(S,\sigma)=\min\left\{k\geq1:
  \begin{array}{l}
   \xi_1,\ldots,\xi_k\in S,\quad \sum_i\xi_i=0,\\[-2pt]
   \prod_i\sigma(\xi_i)=-1
  \end{array}\right\},
 \label{eq:frustration}
\end{equation}
with minimum $\infty$ when the set is empty. Repetitions are allowed.
For an integer $D\geq0$, call the signs \emph{$D$-safe} if $r>D$.
\end{definition}

Every such sign function is $2$-safe: a length-one nonzero frequency
cannot sum to zero, and a length-two relation consists of inverse
frequencies with equal signs.

For a part $A\in\{L,R\}$ write
$\Delta_A(H)=\max_{v\in A}\deg_H(v)$, with maximum zero for an empty
part.

\begin{theorem}[Exact threshold for the termwise method]
\label{thm:threshold}
Fix a real even spectrum with $p>0$, and an integer $D\geq2$. The
following are equivalent:
\begin{enumerate}[label=(\roman*)]
 \item The signs on $S$ are $D$-safe.
 \item Every Fourier-flow weight is nonnegative for every finite simple
 bipartite graph with a designated part $L$ satisfying $\Delta_L(H)\leq D$.
\end{enumerate}
In particular, if
\begin{equation}
 \min(\Delta_L(H),\Delta_R(H))<r(S,\sigma),
 \label{eq:degree-threshold}
\end{equation}
then $t(H,W)\geq p^e$.

If $r=r(S,\sigma)<\infty$, a negative flow already occurs on the graph
$\Theta_{r,3}$ consisting of $r$ internally disjoint three-edge paths
between two endpoints. It has $2r+2$ vertices, $3r$ edges, and maximum
degree $r$ in each part.
\end{theorem}
\begin{proof}
Assume (i), and take a flow of nonzero weight. At a left vertex, discard
all incident zero frequencies; the remaining tuple has length at most
$D$, has sum zero, and has positive sign product. Multiplying these
positive local products over the left part counts every edge exactly
once. Thus the global product has positive sign. A weight containing a
zero coefficient is zero. This proves (ii), and the zero flow proves the
stated Sidorenko inequality.

Conversely, suppose a negative relation of minimum length $r\leq D$
exists, say $\sum_{i=1}^r\xi_i=0$ with sign product $-1$. The endpoints
of $\Theta_{r,3}$ lie in opposite parts. On the $i$th path place the
successive unsigned edge labels
\[
   \xi_i,\ -\xi_i,\ \xi_i.
\]
Every internal vertex has zero sum. Each endpoint has sum
$\sum_i\xi_i=0$. This is a conserved flow, and its weight is
$\prod_i\lambda_{\xi_i}^3$, which is strictly negative. Since $r\geq3$,
all internal degrees are at most $r$, and the graph is simple. Condition
(ii) therefore fails.
\end{proof}

\begin{remark}[What the equivalence does not say]
\label{rem:not-sidorenko}
A negative flow may be outweighed by positive flows. Theorem~
\ref{thm:threshold} classifies the applicability of a \emph{termwise}
proof, not the actual validity of $t(H,W)\geq p^e$. The explicit theta
calculation in Section~\ref{sec:examples} shows that this distinction
persists even for admissible nonnegative hosts.
\end{remark}

\subsection{A finite cutoff and a shortest-witness algorithm}

Let $K=\langle S\rangle\leq\Gamma$, and write
$b(\xi)\in\F_2$ for the sign bit, so that
$\sigma(\xi)=(-1)^{b(\xi)}$. Consider the directed graph on
$K\times\F_2$ whose allowed steps are
\begin{equation}
 (x,a)\longmapsto(x+\xi,a+b(\xi))\qquad(\xi\in S).
 \label{eq:double-cover}
\end{equation}
Because signs are even, every step has its inverse; the graph may also
be viewed as undirected.

\begin{theorem}[Finite detection and algorithm]
\label{thm:bfs}
The frustration length equals the distance from $(0,0)$ to $(0,1)$ in
\eqref{eq:double-cover}, with infinite distance when the target is
unreachable. If it is finite, then
\begin{equation}
 3\leq r(S,\sigma)\leq |K|.
 \label{eq:finite-cutoff}
\end{equation}
Breadth-first search finds a shortest relation, or certifies that none
exists, using $O(|K||S|)$ group operations and $O(|K|)$ stored states,
assuming constant-time group operations. The upper bound $|K|$ is sharp.
\end{theorem}
\begin{proof}
A walk with successive labels $\xi_1,\ldots,\xi_k$ ends at
$(\sum_i\xi_i,\sum_i b(\xi_i))$. Thus reaching $(0,1)$ is exactly a
negative zero-sum relation, including its length.

Take a shortest such relation of length $r$, and let
$s_j=\sum_{i=1}^j\xi_i$ for $0\leq j\leq r$. We claim
$s_0,\ldots,s_{r-1}$ are distinct. Otherwise a nonempty proper
consecutive block is a zero-sum relation. If its sign is negative it
is a shorter witness. If its sign is positive, removing it leaves a
shorter negative zero-sum relation. Both alternatives contradict
minimality. There are only $|K|$ possible partial sums, proving the upper
bound. The lower bound was observed after the definition.

The cover has $2|K|$ states and at most $2|K||S|$ directed edges, so
breadth-first search has the asserted cost and recovers labels from
predecessor pointers. To see sharpness, take $K=\Z/q\Z$ for an odd prime
$q$, $S=\{1,-1\}$, and both signs negative. A shorter zero-sum word has
equal numbers of $1$ and $-1$, hence positive sign, whereas $q$ copies
of $1$ give a negative relation. Thus $r=q=|K|$.
\end{proof}

\begin{corollary}[The sharp Boolean rank cutoff]
\label{cor:boolean-cutoff}
If $K\cong\F_2^k$, then either $r(S,\sigma)=\infty$ or
$r(S,\sigma)\leq k+1$. Hence the terminal relation code is already
reached at degree $k+1$.
\end{corollary}
\begin{proof}
A shortest negative relation cannot repeat a frequency: deleting two
copies preserves its sum and negative sign. Its distinct frequencies
have no nonempty proper zero-sum subset, since either that subset or
its complement would be a shorter negative relation. They form a
minimal linear dependency, so all but one are linearly independent.
Their number is at most $k+1$. Equality of the terminal code follows
by the orthogonal-complement argument used in Theorem~\ref{thm:code}.
Proposition~\ref{prop:circuit} attains the bound for every $k\geq2$.
\end{proof}

The complexity bound is polynomial in the \emph{explicit group size},
not in a compressed presentation of a potentially enormous group. The
algorithm is useful for obtaining a small mathematical certificate even
when listing every zero-sum tuple of a given length would be expensive.

\section{Binary relation codes and global sign characters}
\label{sec:codes}

Choose one representative $\xi_j$ from each inversion orbit in $S$,
for $1\leq j\leq m$. An orbit with $\xi_j=-\xi_j$ is included only once.
For $n=(n_1,\ldots,n_m)\in\Z^m$, put
\[
 \pi(n)=\sum_{j=1}^m n_j\xi_j\in K,
 \qquad \|n\|_1=\sum_j|n_j|.
\]
Write $\bar n\in\F_2^m$ for coordinatewise reduction modulo two.

\begin{definition}[Local relation code]
For $D\geq0$ define
\begin{equation}
 \Rel_D(S)=\Span_{\F_2}
   \{\bar n:\pi(n)=0,\ \|n\|_1\leq D\}
   \ \leq\F_2^m.
 \label{eq:relation-code}
\end{equation}
Write $\Rel_\infty(S)$ for the span with no length restriction.
The safe-sign code is its orthogonal complement
$\Rel_D(S)^\perp$ under the standard dot product on $\F_2^m$.
\end{definition}

\begin{theorem}[Code classification]
\label{thm:code}
Let $b_j\in\F_2$ encode the sign of $\lambda_{\xi_j}$. Then
\begin{equation}
 \sigma\text{ is $D$-safe}
 \quad\Longleftrightarrow\quad
 b\in\Rel_D(S)^\perp.
 \label{eq:safe-code}
\end{equation}
Consequently the number of safe sign assignments is exactly
\begin{equation}
 2^{m-\dim\Rel_D(S)}.
 \label{eq:code-size}
\end{equation}
The spaces $\Rel_D(S)$ increase with $D$, and they stabilize to the full
relation space by $D=|K|$.
\end{theorem}
\begin{proof}
A tuple of frequencies determines an integer vector $n$ by taking the
number of occurrences of $\xi_j$ minus the number of occurrences of
$-\xi_j$. For an orbit of order two one may simply count all occurrences
positively. Its sum is $\pi(n)$, its sign is $(-1)^{b\cdot\bar n}$,
and $\|n\|_1$ is at most the tuple length. The parity identity holds
because subtraction and addition agree modulo two.

Conversely, an integer relation $n$ can be realized by taking $|n_j|$
copies of $\xi_j$ if $n_j\geq0$ and of $-\xi_j$ otherwise. The resulting
tuple has length $\|n\|_1$ and the required parity vector. Thus safety
is precisely orthogonality to all the displayed generators, equivalently
to their span. Rank--nullity proves \eqref{eq:code-size}, and monotonicity
is immediate.

By Theorem~\ref{thm:bfs}, a sign vector satisfying all constraints of
length at most $|K|$ satisfies all constraints of every length. Hence
$\Rel_{|K|}(S)^\perp=\Rel_\infty(S)^\perp$. Taking orthogonal
complements in finite-dimensional $\F_2^m$ proves stabilization.
\end{proof}

\begin{theorem}[The terminal code and translations]
\label{thm:global}
The following conditions are equivalent:
\begin{enumerate}[label=(\roman*)]
 \item $r(S,\sigma)=\infty$.
 \item The sign function extends to a homomorphism $s:K\to\{\pm1\}$.
 \item There is an $a\in G$ such that
       $\chi_\xi(a)=\sigma(\xi)$ for every $\xi\in S$.
 \item Some translate $x\mapsto w(x+a)$ has nonnegative Fourier
       coefficients at every frequency.
\end{enumerate}
The terminal safe-sign code has dimension
\begin{equation}
 \dim_{\F_2}(K/2K).
 \label{eq:terminal-dimension}
\end{equation}
If these conditions hold, $t(H,W)\geq p^e$ for every bipartite $H$.
\end{theorem}
\begin{proof}
Under (i), every relation among the generators has positive sign, so
\[
 s\left(\sum_j n_j\xi_j\right)=\prod_j\sigma(\xi_j)^{n_j}
\]
is well defined and is a homomorphism. The converse is immediate by
multiplying the signs along a zero-sum relation.

Every circle-valued character of a subgroup of a finite abelian group
extends to the whole group. For completeness, suppose it is defined on
$K_0$, and adjoin $\gamma$ whose coset has order $n$ modulo $K_0$.
Choose $z\in\T$ with $z^n=s(n\gamma)$ and set
$s'(k+j\gamma)=s(k)z^j$. The choice of root makes this well defined.
Iterating proves extension to $\Gamma$. Finite Fourier duality identifies
characters of $\Gamma$ with evaluations $\xi\mapsto\chi_\xi(a)$ for
$a\in G$. This proves (ii)$\Rightarrow$(iii); the reverse follows by
restriction, because the values on the generating set are signs.

The Fourier coefficient of $w(x+a)$ at $\xi$ is
$\lambda_\xi\chi_\xi(a)$. It is nonnegative for every active frequency
exactly when the multiplier is $\sigma(\xi)$, proving the equivalence
with (iv). Homomorphisms $K\to\{\pm1\}$ are precisely the linear
functionals on $K/2K$, and distinct ones differ on $S$. This gives
\eqref{eq:terminal-dimension}. Finally, Theorem~\ref{thm:threshold}
applies with arbitrarily large $D$.
\end{proof}

\begin{remark}[The translating element need not have order two]
It is not valid to require the sign character to extend as a
$\{\pm1\}$-valued character on all of $\Gamma$. For example, the
nontrivial sign character on $2\Z/4\Z$ cannot do so. It extends as a
circle-valued character, corresponding to a translate of order four.
The correct assertion is (iii), only on the active spectrum. In that
case the translated function is still real and even, because its active
Fourier coefficients are real and invariant under inversion.
\end{remark}

There is also a direct density interpretation: replacing $w(x-y)$ by
$w(x-y+a)$ does not change a bipartite homomorphism density, since one
can translate all variables in $L$ by $a$. The group-theoretic
classification therefore recovers global positivity after a suitable
translation. The finite-degree codes are substantially larger, as the
next section shows.

\section{Strict hierarchies and a cautionary theta calculation}
\label{sec:examples}

\subsection{Independent frequencies in odd characteristic}

\begin{proposition}[A dimension-independent degree window]
\label{prop:odd-spectrum}
Let $G=\F_q^m$, where $q$ is an odd prime, and let
$S=\{\pm e_1,\ldots,\pm e_m\}$ in the dual group. Every even assignment
of signs on $S$ is $(q-1)$-safe. If any sign is negative, its frustration
length is exactly $q$, and no translate makes every Fourier coefficient
nonnegative.
\end{proposition}
\begin{proof}
An integer relation of $\ell^1$ length less than $q$ has each coordinate
congruent to zero modulo $q$ and of absolute value less than $q$. Every
coordinate is therefore zero as an integer, so its parity vector is zero.
All signs are safe through degree $q-1$. If the sign of $e_j$ is negative,
$q$ copies of $e_j$ give a negative zero-sum relation, so the threshold is
exactly $q$. Finally $K=\F_q^m$ has odd order, and consequently has no
nontrivial homomorphism to $\{\pm1\}$. Theorem~\ref{thm:global} proves
the translation assertion.
\end{proof}

For example, take
\begin{equation}
 w(x)=p+2\sum_{j=1}^m a_j\cos(2\pi x_j/q),
 \qquad
 2\sum_j|a_j|\leq\min(p,1-p),
 \label{eq:odd-example}
\end{equation}
with all $a_j\ne0$. Some $a_j$ may be negative. This is an admissible
host that satisfies the Sidorenko inequality for \emph{every} bipartite
pattern whose smaller one-sided maximum degree is less than $q$.
There is no restriction on $m$. Thus the criterion is not simply global
nonnegative Fourier spectrum in disguise.

\subsection{Every finite threshold in characteristic two}

\begin{proposition}[Circuit spectra]
\label{prop:circuit}
For every integer $r\geq3$ there is a weighted scalar Cayley host on
$\F_2^{r-1}$ whose frustration length is exactly $r$.
\end{proposition}
\begin{proof}
Take
\[
 S=\{e_1,\ldots,e_{r-1},e_1+\cdots+e_{r-1}\}.
\]
The only nonzero linear dependency among these $r$ distinct frequencies
is their full sum. Assign signs with total product $-1$. A zero-sum
integer tuple has, modulo two, either no odd multiplicities or odd
multiplicity at all $r$ frequencies. The former has positive sign; the
latter has length at least $r$. The tuple containing all $r$ frequencies
once is a negative relation. Lemma~\ref{lem:realizable} supplies an
admissible $w$ with precisely these signs.
\end{proof}

In particular, the descending safe-sign spaces can lose a dimension
at an arbitrarily late prescribed degree. The phenomenon is not an
artifact of using a cyclic group of large odd order.

\subsection{Graphical spectra and signed cycles}

There is a useful concrete coding interpretation. Let $X$ be a finite
simple graph, and attach the frequency $e_u+e_v\in\F_2^{V(X)}$ to each
edge $uv$ of $X$. Assign signs to those frequencies.

\begin{proposition}[A cycle-space realization]
\label{prop:graphical}
For the graphical spectrum just described, the frustration length is
the length of the shortest cycle in $X$ with negative edge-sign product,
or infinity if no such cycle exists. Moreover $\Rel_D(S)$ is the span
of the edge-incidence vectors of cycles of length at most $D$.
\end{proposition}
\begin{proof}
Reducing a frequency tuple modulo two gives an edge set in which every
vertex of $X$ has even degree. Its sign equals the sign of the tuple,
and its cardinality is no greater than the tuple length. Every nonempty
even-degree edge set decomposes into edge-disjoint simple cycles: follow
unused edges until returning to a previously visited vertex, remove the
resulting cycle, and iterate. A negative total sign forces at least one
negative cycle. This proves the threshold assertion in both directions.

For the code assertion, an integer relation of length at most $D$ reduces
to an even-degree edge set with at most $D$ edges. Every cycle in its
edge-disjoint decomposition has length at most $D$. Conversely, each
such cycle itself gives an integer relation of that length. Hence the
spans agree.
\end{proof}

This is a realization of ordinary cycle-space reasoning inside the
Fourier obstruction code, not a claim of a new theory of signed graphs.
It provides easily inspectable examples and compact parity certificates.

\subsection{Why a negative term is not a counterexample}

The theta witness in Theorem~\ref{thm:threshold} is optimal for the
\emph{sign test}, but it need not violate Sidorenko. This can be seen
without appealing to any general theorem about theta graphs.

Let $q$ be an odd prime, $G=\Z/q\Z$, and
\begin{equation}
 w(x)=p+a\chi(x)+a\overline{\chi(x)},
 \qquad a<0,
 \qquad 2|a|\leq\min(p,1-p),
 \label{eq:theta-w}
\end{equation}
where $\chi(x)=\exp(2\pi ix/q)$. The only nonzero Fourier coefficients
are $\lambda_0=p$ and $\lambda_{\pm1}=a$.

\begin{proposition}[Exact theta polynomial]
\label{prop:theta-polynomial}
For $\Theta_{q,3}$ and \eqref{eq:theta-w},
\begin{equation}
 t(\Theta_{q,3},W)=
 \sum_{j=0}^{(q-1)/2}
   \binom q{2j}\binom{2j}{j}
   p^{3(q-2j)}a^{6j}
   +2a^{3q}.
 \label{eq:theta-poly}
\end{equation}
The last term is negative, but $t(\Theta_{q,3},W)\geq p^{3q}$.
\end{proposition}
\begin{proof}
For normalized convolution
$(f*g)(z)=\E_y f(z-y)g(y)$, averaging the two internal vertices of a
three-edge path gives $h=w*w*w$. Thus
$t(\Theta_{q,3},W)=\E_z h(z)^q$, and
\[
 h(z)=p^3+a^3\chi(z)+a^3\overline{\chi(z)}.
\]
In expanding its $q$th power, a term with $u$ positive and $v$ negative
frequencies survives averaging exactly when $u-v\equiv0\pmod q$.
Since $u+v\leq q$, either $u=v=j$ or $(u,v)=(q,0),(0,q)$.
These alternatives give exactly \eqref{eq:theta-poly}.

On the other hand, $w\geq0$ implies $h\geq0$, and $\E h=p^3$.
Convexity of $x\mapsto x^q$ on $[0,\infty)$ gives
$\E h^q\geq(\E h)^q=p^{3q}$. The negative Fourier contribution is
therefore genuinely present but is compensated by the others.
\end{proof}

\section{Moment defects and the robust cycle-inventory inequality}
\label{sec:robust}

The exact sign test is qualitative. We now measure its failure without
requiring a spectral support cutoff or discarding small coefficients.
Define the real even Fourier-absolute companion
\begin{equation}
 w^\#(x)=\sum_{\xi\in\Gamma}|\lambda_\xi|\chi_\xi(x).
 \label{eq:abs-companion}
\end{equation}
There is no assertion that $w^\#$ takes values in $[0,1]$, or even that
it is pointwise nonnegative. Its usefulness comes from its Fourier
coefficients.

For every integer $d\geq0$ set
\begin{align}
 M_d&=\E_x w^\#(x)^d
   =\sum_{\xi_1+\cdots+\xi_d=0}
                  \prod_{i=1}^d|\lambda_{\xi_i}|,\label{eq:M}\\
 B_d&=\E_x w(x)^d
   =\sum_{\xi_1+\cdots+\xi_d=0}
                  \prod_{i=1}^d\lambda_{\xi_i},\label{eq:B}\\
 D_d&=\frac{M_d-B_d}{2}.\label{eq:D}
\end{align}
Empty products give $M_0=B_0=1$ and $D_0=0$.

\begin{lemma}[Exact local negative mass]
\label{lem:defect}
The number $D_d$ is the sum of the absolute weights of all zero-sum
frequency $d$-tuples whose coefficient product is negative. In particular
$D_d\geq0$ and $D_0=D_1=D_2=0$. If $w\geq0$, then
\begin{equation}
 0<p^d\leq B_d\leq M_d\qquad(d\geq1).
 \label{eq:moments-positive}
\end{equation}
For any $D\geq0$, the signs are $D$-safe if and only if
$D_d=0$ for all $0\leq d\leq D$; equivalently, it is enough that
$D_D=0$.
\end{lemma}
\begin{proof}
Orthogonality proves the moment identities. Subtracting the signed sum
from its absolute-weight sum cancels each nonnegative term and doubles
the absolute value of each negative term. This proves the first claim.
There are no negative relations of lengths one or two; zero frequencies
have positive coefficient. For nonnegative $w$, Jensen's inequality
proves $B_d\geq p^d$ for $d\geq1$, with equality for $d=1$.
Finally, every negative active tuple has strictly positive absolute
weight. Removing its zero labels leaves precisely a relation of the
type in \eqref{eq:frustration}. These observations prove both directions
of the last assertion. If a negative relation has length $k\leq D$,
padding it by $D-k$ zero frequencies multiplies its negative weight by
the positive factor $p^{D-k}$. Thus $D_D=0$ alone is sufficient.
\end{proof}

For a part $A=L$ or $R$ define
\begin{equation}
 \Err_A(H,w)=\frac12\left(
   \prod_{v\in A}M_{\deg(v)}-
   \prod_{v\in A}B_{\deg(v)}\right),
 \qquad
 \Err_H(w)=\min\{\Err_L(H,w),\Err_R(H,w)\}.
 \label{eq:error}
\end{equation}
These expressions are nonnegative, including for signed $w$: they will
be identified as negative masses in a partially constrained sum.

Let $c_\ell(H)$ count undirected simple cycles of length $\ell$ in $H$,
each cycle counted once, not once per starting point or orientation.
Define
\begin{equation}
 Q_\ell(w)=\sum_{\xi\ne0}|\lambda_\xi|^\ell,
 \qquad
 \mathcal C_H(w)=\sum_{\substack{\ell\geq4\\\ell\text{ even}}}
       c_\ell(H)p^{e-\ell}Q_\ell(w).
 \label{eq:inventory}
\end{equation}
Only lengths of cycles actually present in $H$ are included, so every
power $p^{e-\ell}$ in this sum has nonnegative exponent.

\begin{theorem}[Robust cycle inventory]
\label{thm:inventory}
For every finite simple bipartite $H$ and every real even $w$ with
$p>0$,
\begin{equation}
 \boxed{\quad
 t(H,W)\ \geq\ p^e+\mathcal C_H(w)-\Err_H(w).
 \quad}
 \label{eq:main-inventory}
\end{equation}
If all full Fourier-flow weights are nonnegative, the error in
\eqref{eq:main-inventory} can instead be taken to be zero. In particular,
zero error is valid under the one-sided sign condition
\eqref{eq:degree-threshold}.
\end{theorem}
\begin{proof}
Let $N_H$ be the absolute mass of the negative terms in
\eqref{eq:flow-expansion}. Fix one side $A$. Relax the flow constraints
by imposing them only at vertices of $A$. Each edge is incident to
exactly one such vertex, so the variables at distinct vertices of $A$
are disjoint. The total absolute weight of the relaxed assignments is
$\prod_{v\in A}M_{\deg(v)}$, whereas their signed total is
$\prod_{v\in A}B_{\deg(v)}$. Their negative mass is therefore exactly
$\Err_A(H,w)$. Every negative full flow is among these assignments,
with exactly the same absolute weight. Thus
\begin{equation}
 N_H\leq\Err_A(H,w),\qquad N_H\leq\Err_H(w).
 \label{eq:negative-control}
\end{equation}

The positive-weight sum contains the zero flow, of weight $p^e$. It
also contains the following flows. For each simple cycle $C$ of even
length $\ell$ and each $\xi\ne0$, put alternating labels
$\xi,-\xi$ around $C$ and zero labels elsewhere. Choose one distinguished
edge of each cycle to carry $\xi$, thereby fixing the alternating rule.
This is a flow of weight
$p^{e-\ell}|\lambda_\xi|^\ell\geq0$. If the coefficient is nonzero,
its edge support is exactly $C$. Distinct cycles therefore produce
distinct flows; for a fixed cycle the distinguished-edge label separates
the choices of $\xi$. Frequencies of order two cause no duplication,
since such a frequency appears only once in the sum over $\Gamma$.
Zero-weight assignments may harmlessly be omitted.

It follows that the positive mass is at least
$p^e+\mathcal C_H(w)$. Subtracting the negative mass and using
\eqref{eq:negative-control} proves the result. Notice the factor $1/2$
in \eqref{eq:error}: it is essential because we are bounding the
negative mass, not the difference between the original expansion and an
all-absolute expansion.
\end{proof}

The estimate retains only a subset of positive flows. More complicated
Eulerian supports can make it strict. Its advantage is that the retained
terms have exact coefficients and require no classification of all flows.

\subsection{Relative local defects and inverse certificates}

Assume $w\geq0$. Write
\begin{equation}
 \eps_d=\frac{D_d}{M_d},\qquad
 P_A=\prod_{v\in A}M_{\deg(v)},
 \label{eq:relative-defect}
\end{equation}
with $\eps_0=0$. Lemma~\ref{lem:defect} gives
$0\leq\eps_d<1/2$. Equation~\eqref{eq:error} becomes the exact identity
\begin{equation}
 \Err_A=\frac{P_A}{2}
   \left(1-\prod_{v\in A}(1-2\eps_{\deg(v)})\right)
 \leq P_A\sum_{v\in A}\eps_{\deg(v)}.
 \label{eq:relative-error}
\end{equation}
The inequality follows by expanding, or inductively from
$1-\prod_i(1-x_i)\leq\sum_i x_i$ for $x_i\in[0,1]$.

\begin{corollary}[A deficit forces local negative mass]
\label{cor:inverse}
Suppose $w\geq0$ and
$\eta=p^e-t(H,W)>0$. Then on \emph{both} sides $A=L,R$,
\begin{equation}
 \Err_A(H,w)\geq\eta+\mathcal C_H(w).
 \label{eq:inverse-error}
\end{equation}
Consequently, for each side there is a vertex $v\in A$ with
\begin{equation}
 \frac{D_{\deg(v)}}{M_{\deg(v)}}
 \geq\frac{\eta+\mathcal C_H(w)}{|A|P_A}.
 \label{eq:inverse-vertex}
\end{equation}
In particular, the spectrum has a sign-frustrated zero-sum relation of
length at most $\deg(v)$.
\end{corollary}
\begin{proof}
Rearrange \eqref{eq:main-inventory}; since $\Err_H$ is the minimum of
the two side errors, each side is at least the displayed quantity.
Equation~\eqref{eq:relative-error} and averaging over $A$ prove
\eqref{eq:inverse-vertex}. A positive defect contains a negative tuple;
discarding its zero labels proves the final assertion. A positive deficit
cannot occur when either side is empty, so the denominators are valid.
\end{proof}

The estimate is independent of $|G|$ in form, but its moment factors can
be large. It is not a bound solely in terms of the density $p$ or the
cardinality of the spectrum. Any simplification of those factors requires
additional analytic or structural hypotheses.

\subsection{The $U^2$ interpretation}

For any complex function $f:G\to\C$, use the normalized convention
\begin{equation}
 \|f\|_{U^2}^4=
 \E_{x,a,b}f(x)\overline{f(x+a)}\,
                 \overline{f(x+b)}f(x+a+b).
 \label{eq:u2}
\end{equation}
Expanding the four Fourier series and averaging successively in $x,a,b$
forces the same frequency at all four positions. Hence
\begin{equation}
 \|f\|_{U^2}^4=\sum_{\xi\in\Gamma}|\widehat f(\xi)|^4.
 \label{eq:u2-fourier}
\end{equation}
For $f=w-p$, its zero coefficient vanishes, so
\begin{equation}
 Q_4(w)=\|w-p\|_{U^2}^4.
 \label{eq:Q4-U2}
\end{equation}
This is the normalization used in all subsequent bounds.

For $G=\Z/N\Z$, the unnormalized Fourier transform in the inspected
\texttt{ProveIt} definitions must be divided by $N$ to match
\eqref{eq:fourier}; reversing the frequency convention has no effect
for an even real $w$. The raw two-dimensional cube sum has $N^3$
terms. Thus the corresponding unnormalized fourth-root quantity is
$N^{3/4}\|f\|_{U^2}$, not $\|f\|_{U^2}$ itself. This conversion is
important for any Lean integration \cite{ProveItDefinitions}.

\section{Sharp girth forcing}
\label{sec:forcing}

For a real kernel on $G\times G$, use the set-cut convention
\begin{equation}
 \Cut{K}=\max_{A,B\subseteq G}
    \left|\E_{x,y}\one_A(x)K(x,y)\one_B(y)\right|.
 \label{eq:cut}
\end{equation}
Different conventions allowing arbitrary $[-1,1]$ test functions have
different constants; the factor $1/4$ below is for \eqref{eq:cut}.
Write
\begin{equation}
 \beta(w)=\max_{\xi\ne0}|\lambda_\xi|,
 \label{eq:beta}
\end{equation}
with maximum zero when there is no nonprincipal frequency.

\begin{lemma}[The cut-norm constant]
\label{lem:cut}
For $W(x,y)=w(x-y)$ and $p=\E w$,
\begin{equation}
 \Cut{W-p}\leq\frac{\beta(w)}4.
 \label{eq:cut-beta}
\end{equation}
\end{lemma}
\begin{proof}
The convolution operator
$Tf(x)=\E_y(w(x-y)-p)f(y)$ kills constants and acts on nonprincipal
characters with eigenvalues $\lambda_\xi$. Parseval gives its $L^2$
operator norm as $\beta(w)$. For sets $A,B$, subtract their densities
$\alpha,\gamma$ from their indicators. The bilinear form is unchanged,
and Cauchy--Schwarz bounds it by
\[
 \beta(w)\sqrt{\alpha(1-\alpha)}
              \sqrt{\gamma(1-\gamma)}\leq\beta(w)/4.
\]
Taking the maximum proves the claim.
\end{proof}

\begin{theorem}[Robust forcing with the girth exponent]
\label{thm:forcing}
Let $H$ be a finite simple bipartite graph containing a cycle, let its
girth be $g$, and let $c_g=c_g(H)>0$. Let $0\leq w\leq1$ be even,
with $p>0$. Put
\begin{equation}
 \delta=\frac{t(H,W)}{p^e}-1,
 \qquad \eps=\frac{\Err_H(w)}{p^e}.
 \label{eq:global-delta}
\end{equation}
Then $\delta+\eps\geq0$, and
\begin{align}
 Q_g(w)&\leq\frac{p^g}{c_g}(\delta+\eps),\label{eq:force-Q}\\
 \beta(w)&\leq\frac{p}{c_g^{1/g}}(\delta+\eps)^{1/g},\label{eq:force-beta}\\
 \Cut{W-p}&\leq\frac{p}{4c_g^{1/g}}
                    (\delta+\eps)^{1/g}.
 \label{eq:force-cut}
\end{align}
Whenever all full Fourier-flow weights are nonnegative, one can use
$\eps=0$ instead, whether or not the particular local error functional
in \eqref{eq:error} vanishes. For $g=4$, \eqref{eq:force-Q} also gives
\begin{equation}
 \|w-p\|_{U^2}\leq\frac{p}{c_4^{1/4}}(\delta+\eps)^{1/4}.
 \label{eq:force-u2}
\end{equation}
\end{theorem}
\begin{proof}
Retain only the shortest-cycle term from Theorem~\ref{thm:inventory}:
\[
 c_gp^{e-g}Q_g(w)\leq t(H,W)-p^e+\Err_H(w).
\]
Divide by $p^e$ to obtain \eqref{eq:force-Q}. The inequality
$\beta(w)^g\leq Q_g(w)$ gives \eqref{eq:force-beta};
Lemma~\ref{lem:cut} gives \eqref{eq:force-cut}. The zero-error version
follows from the zero-error clause of Theorem~\ref{thm:inventory}.
Equation~\eqref{eq:Q4-U2} proves the final assertion.
\end{proof}

\subsection{Both exponent and leading constant are optimal}

\begin{theorem}[Sharpness in the weighted scalar class]
\label{thm:sharp}
For every fixed cyclic bipartite $H$ of girth $g$, and every fixed
$p\in(0,1)$, the exponent $1/g$ in \eqref{eq:force-cut} cannot be
replaced by a larger exponent in a uniform zero-error bound. With
exponent $1/g$, the leading constant $1/(4c_g^{1/g})$ cannot be
reduced. The coefficient $c_g$ in \eqref{eq:force-Q} is likewise
optimal.
\end{theorem}
\begin{proof}
Take $G=\Z/2\Z$, with its nontrivial character $\chi$, and let
\begin{equation}
 w_t(x)=p+t\chi(x),\qquad
 0<|t|\leq\min(p,1-p).
 \label{eq:rank-one}
\end{equation}
This is an admissible host. A conserved labeling selects the edges
carrying the nonzero frequency. The selected edge set $F$ must have
even degree at every vertex. Conversely every such $F$ is a conserved
labeling. Therefore
\begin{equation}
 t(H,W_t)=\sum_{\substack{F\subseteq E\\
                         \deg_F(v)\ \mathrm{even}\ \forall v}}
                  p^{e-|F|}t^{|F|}.
 \label{eq:even-subgraph}
\end{equation}
Every even-degree edge set decomposes into cycles. Since $H$ is
bipartite, $|F|$ is even, so every term is nonnegative even when $t<0$.
The smallest nonempty such sets have size $g$ and are exactly the
$c_g$ shortest cycles. The next possible size is at least $g+2$.
Consequently
\begin{equation}
 \delta_t:=\frac{t(H,W_t)}{p^e}-1
   =c_g(t/p)^g+O_H\bigl((t/p)^{g+2}\bigr).
 \label{eq:sharp-asymptotic}
\end{equation}
The meaning of the remainder is elementary: a finite polynomial has
only powers at least $g+2$ left after its leading term. If no such powers
occur, the remainder is identically zero.

Here $\beta(w_t)=|t|$ and $Q_g(w_t)=|t|^g$. Moreover
$W_t(x,y)-p=t\chi(x)\chi(y)$. Taking both sets in \eqref{eq:cut} to
be the $\chi=1$ class gives $|t|/4$, so
\begin{equation}
 \Cut{W_t-p}=|t|/4.
 \label{eq:exact-rank-one-cut}
\end{equation}
It follows that
\begin{equation}
 \lim_{t\to0}
 \frac{\Cut{W_t-p}}{p\delta_t^{1/g}}
      =\frac1{4c_g^{1/g}}.
 \label{eq:sharp-limit}
\end{equation}
Any larger exponent would make the proposed right-hand side decay
strictly faster than $|t|$, and any smaller constant with the same
exponent would contradict this limit. The corresponding limit for
$Q_g$ proves its coefficient optimal. These hosts also have zero local
sign defect, since a zero-sum tuple over $\Z/2\Z$ uses the nonzero
frequency an even number of times.
\end{proof}

Sharpness here permits weighted kernels and diagonal values. It does
not by itself prove a sharp constant in the narrower class of loopless
$0$--$1$ Cayley graphs; that transfer is a separate question.

\subsection{The indispensable forest exception}

\begin{proposition}[Forests do not force in regular hosts]
\label{prop:forest}
For every forest $H$ and every scalar Cayley host of mean $p$,
\begin{equation}
 t(H,W)=p^{e(H)}.
 \label{eq:forest}
\end{equation}
No conclusion $\Cut{W-p}\to0$ follows from this equality alone.
\end{proposition}
\begin{proof}
An isolated vertex contributes one. If an edge is present, integrate a
leaf variable first; its only incident factor averages to $p$,
independently of the remaining endpoint. Delete the leaf and iterate.
The nonconstant family \eqref{eq:rank-one} has positive cut norm but
still satisfies \eqref{eq:forest}, proving the last statement.
\end{proof}

For the unnormalized weighted adjacency matrix
$A_{xy}=w(x-y)$, the eigenvalues are $|G|\lambda_\xi$, whereas the
normalized convolution operator has eigenvalues $\lambda_\xi$.
For example, a bound $\beta(w)\leq p\rho$ means nonprincipal adjacency
eigenvalues of size at most $|G|p\rho$, not $p\rho$. This normalization
and the cyclic-pattern hypothesis are essential when comparing forcing
statements such as the concluding discussion in \cite{YuqiZhao}.

\subsection{Even subdivisions as an established positivity class}

\begin{proposition}[Fourier pairing on even subdivisions]
\label{prop:subdivision}
Suppose $H$ is obtained by replacing every edge of a graph by a path
of positive even length. For every real even $w$, every full Fourier-flow
weight of $H$ is nonnegative. Hence, when $p>0$, the zero-error
cycle-inventory inequality holds; when $H$ is cyclic, so do the sharp
girth-forcing bounds.
\end{proposition}
\begin{proof}
The original vertices lie in one part of the subdivided graph. On a
replacement path, the degree-two conservation equations force
successive labels to be $\xi,-\xi,\xi,-\xi,\ldots$. A path of even
length $2k$ therefore contributes $\lambda_\xi^{2k}\geq0$. Multiplying
over all replacement paths gives a nonnegative global weight. The
remaining conclusions follow from Theorems~\ref{thm:inventory} and
\ref{thm:forcing}.
\end{proof}

The baseline Sidorenko conclusion in this proposition is not presented
as new; it is the positivity mechanism in \cite{YuqiZhao}. The useful
additional point here is the sharp shortest-cycle forcing statement.
The particular one-sided moment defect $\Err_H$ need not vanish for
this argument to work: positivity of complete flows is enough.

\section{Block-local certification}
\label{sec:blocks}

High degrees at articulation vertices can make a global maximum-degree
test unnecessarily restrictive. Scalar translation invariance permits
an exact reduction to blocks.

\begin{lemma}[Gluing at one vertex]
\label{lem:glue}
Suppose $H$ is the union of graphs $H_1,H_2$ that meet at exactly one
vertex and have disjoint edge sets. Then for every scalar Cayley host,
\begin{equation}
 t(H,W)=t(H_1,W)t(H_2,W).
 \label{eq:gluing}
\end{equation}
\end{lemma}
\begin{proof}
Fix the image $x$ of the common vertex. The remaining variables in the
two graphs are independent. The rooted homomorphism density of either
graph is independent of $x$, because translating every variable by
$-x$ leaves every difference unchanged. Its constant value equals the
unrooted density. Multiplication and averaging over $x$ prove the claim.
\end{proof}

Treat bridges as single-edge blocks, and let $\mathcal B$ denote the
blocks containing cycles. Repeated use of \eqref{eq:gluing}, component
factorization, and leaf integration yields
\begin{equation}
 t(H,W)=p^{b(H)}\prod_{B\in\mathcal B}t(B,W),
 \label{eq:block-factor}
\end{equation}
where $b(H)$ is the number of bridges. Isolated vertices contribute one.

\begin{corollary}[A blockwise degree threshold]
\label{cor:block-degree}
Let $w$ be even with $p>0$. If
\begin{equation}
 r(S,\sigma)>
 \max_{B\in\mathcal B}
       \min\{\Delta_{L\cap V(B)}(B),
              \Delta_{R\cap V(B)}(B)\},
 \label{eq:block-degree}
\end{equation}
then $t(H,W)\geq p^{e(H)}$. The side chosen in the minimum may vary
from block to block.
\end{corollary}
\begin{proof}
Apply Theorem~\ref{thm:threshold} to each cyclic block and multiply the
resulting positive lower bounds in \eqref{eq:block-factor}.
\end{proof}

For a robust version assume $w\geq0$, and put
$q_B=\min\{1,\Err_B(w)/p^{e(B)}\}$. Nonnegativity of every block density
and Theorem~\ref{thm:inventory}, with its surplus discarded, give
\begin{equation}
 \frac{t(H,W)}{p^{e(H)}}\geq\prod_{B\in\mathcal B}(1-q_B).
 \label{eq:robust-block}
\end{equation}
This does not replace the sharper full-graph inventory in all examples;
it is an alternative certificate when block degrees or block defects
are substantially smaller. The exact factorization is a consequence
of translation invariance, not a generic statement about arbitrary
irregular graphons.

\section{The 35-vertex incidence pattern}
\label{sec:incidence}

\subsection{The finite object and its short cycles}

Let $I=\{0,\ldots,12\}$ and let $J=\{0,\ldots,21\}$ index the triples
in Table~\ref{tab:faces}. Define $\Hstar$ to be their incidence graph,
with disjoint vertex sets $I$ and $J$, and with $i$ adjacent to $j$ exactly
when point $i$ occurs in face $j$. This is the finite pattern specified
in \cite{OpenAIComplex}; we use its explicit data, not an assertion about
a host witnessing its failure.

\begin{table}[htbp]
\centering
\begin{tabular}{r@{\quad}l@{\qquad\qquad}r@{\quad}l}
\toprule
Face & Points & Face & Points\\
\midrule
0 & $0,1,3$ & 11 & $3,10,11$\\
1 & $0,1,9$ & 12 & $4,8,9$\\
2 & $0,2,3$ & 13 & $4,8,11$\\
3 & $0,2,9$ & 14 & $5,6,7$\\
4 & $1,3,10$ & 15 & $5,6,11$\\
5 & $1,7,10$ & 16 & $5,7,10$\\
6 & $1,7,12$ & 17 & $5,10,11$\\
7 & $1,9,12$ & 18 & $6,7,12$\\
8 & $2,3,4$ & 19 & $6,8,11$\\
9 & $2,4,9$ & 20 & $6,8,12$\\
10 & $3,4,11$ & 21 & $8,9,12$\\
\bottomrule
\end{tabular}
\caption{The complete face list defining $\Hstar$. Numerical incidence
data are from the finite-complex section of \cite{OpenAIComplex}.}
\label{tab:faces}
\end{table}

\begin{proposition}[Short-cycle certificate]
\label{prop:incidence-data}
The graph $\Hstar$ is connected, with 35 vertices and 66 edges. Every
face vertex has degree three; the point degrees in order are
\begin{equation}
 (4,6,4,6,5,4,5,5,5,6,5,6,5).
 \label{eq:point-degrees}
\end{equation}
Its girth is four, and
\begin{equation}
 c_4(\Hstar)=33,\qquad c_6(\Hstar)=88.
 \label{eq:incidence-cycles}
\end{equation}
\end{proposition}
\begin{proof}
There are 22 triples on 13 points, giving 35 vertices and 66 incidences.
Counting occurrences yields \eqref{eq:point-degrees}. In the displayed
face order, every face after the first meets an earlier face, and all
13 points occur; hence the incidence graph is connected.

List the point pairs occurring in the triples. There are 33 distinct
pairs, each occurring in exactly two faces; this can be checked directly
from the complete list. A four-cycle is uniquely a choice of a point
pair and two faces containing it. Consequently
$c_4=\sum_{\{i,k\}}\binom{2}{2}=33$. Simplicity and bipartiteness
exclude shorter cycles.

For the six-cycle count, form the point graph whose edges are those 33
pairs. Its three-cliques, checked from the same list, are exactly the
22 displayed faces. Any six-cycle in the incidence graph uses three
points that form such a clique. For a fixed face $\{a,b,c\}$, each of
its three point pairs has two choices of containing face. All three
pair-face sets have the given face in common, and any two intersect
only there: a face containing two of these pairs contains the entire
triple. Among the $2^3=8$ assignments, the assignments with three
distinct faces number
$8-3\cdot2+3\cdot1-1=4$. These are precisely the six-cycles on the
given point triple. The total is $22\cdot4=88$.

The distributed verifier separately checks pair multiplicities, the
three-clique list, connectedness, and these counts using exact integers.
\end{proof}

\subsection{Cubic defects control all scalar-host violations}

Put $f=w-p$, and abbreviate
\begin{equation}
 E_3(w)=\frac{M_3^{22}-B_3^{22}}2,
 \qquad
 R_*(w)=33p^{62}\|f\|_{U^2}^4+88p^{60}Q_6(w).
 \label{eq:incidence-error}
\end{equation}

\begin{theorem}[Incidence-pattern inequality and inverse obstruction]
\label{thm:incidence}
Let $G$ be any finite abelian group and let $0\leq w\leq1$ be even with
$p=\E w>0$. Then
\begin{equation}
 \boxed{\quad t(\Hstar,W)\geq p^{66}+R_*(w)-E_3(w).\quad}
 \label{eq:incidence-main}
\end{equation}
If $D_3=0$, the error vanishes, and in particular
\begin{equation}
 t(\Hstar,W)\geq p^{66}+33p^{62}\|w-p\|_{U^2}^4.
 \label{eq:incidence-u2}
\end{equation}
In that case, with $\delta=t(\Hstar,W)/p^{66}-1$,
\begin{equation}
 \|w-p\|_{U^2}\leq p(\delta/33)^{1/4},\qquad
 \Cut{W-p}\leq\frac{p}{4\cdot33^{1/4}}\delta^{1/4}.
 \label{eq:incidence-force}
\end{equation}
Both constants and the exponent in \eqref{eq:incidence-force} are
asymptotically sharp in this class.

Conversely, if $\eta=p^{66}-t(\Hstar,W)>0$, then
\begin{equation}
 \boxed{\quad
 D_3\geq\frac{\eta+R_*(w)}{22M_3^{21}}>0.
 \quad}
 \label{eq:incidence-inverse}
\end{equation}
Thus a scalar Cayley counterexample for this pattern must have
$\xi+\zeta+\theta=0$ with
$\lambda_\xi\lambda_\zeta\lambda_\theta<0$, and the entire collection
of such triples has at least the weighted mass in
\eqref{eq:incidence-inverse}.
\end{theorem}
\begin{proof}
The face part consists of 22 degree-three vertices. Its error in
\eqref{eq:error} is exactly $E_3(w)$, so $\Err_{\Hstar}(w)\leq E_3(w)$.
Retain the four- and six-cycle terms of Theorem~\ref{thm:inventory},
use Proposition~\ref{prop:incidence-data}, and apply
\eqref{eq:Q4-U2}. This gives \eqref{eq:incidence-main}.

If $D_3=0$, then $M_3=B_3$, so the error is zero. A relation of length
one or two cannot be negative, and a negative triple containing a zero
frequency would reduce to one of those cases. Thus $D_3=0$ is also
exactly three-safety. The $U^2$ and cut bounds follow from the preceding
section, with $g=4$ and $c_g=33$; their sharpness follows from
Theorem~\ref{thm:sharp}.

Finally a deficit implies $E_3(w)\geq\eta+R_*(w)$. Since
$0<B_3\leq M_3$ and $M_3-B_3=2D_3$,
\[
 E_3(w)=D_3\sum_{j=0}^{21}M_3^{21-j}B_3^j
       \leq22D_3M_3^{21}.
\]
This proves \eqref{eq:incidence-inverse}. Lemma~\ref{lem:defect}
identifies $D_3$ with the claimed exact weighted triple mass.
\end{proof}

A useful normalized robust form, even when $D_3\ne0$, is obtained by
putting $\eps_*=E_3(w)/p^{66}$:
\begin{equation}
 \|w-p\|_{U^2}\leq p\bigl((\delta+\eps_*)/33\bigr)^{1/4},
 \qquad
 \Cut{W-p}\leq\frac p4
       \bigl((\delta+\eps_*)/33\bigr)^{1/4}.
 \label{eq:incidence-robust-force}
\end{equation}
Here $\delta$ may be negative, but the sum in the fourth root is
nonnegative by \eqref{eq:incidence-main}.

The theorem gives a restricted-host forcing result for a pattern used
in a general nonforcing claim. There is no conflict: a host with cubic
sign frustration lies outside the zero-defect class, and a typed
activation kernel is not automatically a scalar convolution kernel.
Whether this particular pattern has a scalar abelian counterexample
remains a question, not a consequence of either construction.

\section{A sufficient certificate for typed convolution kernels}
\label{sec:typed}

The source construction motivates keeping vertex types rather than
silently replacing them by a scalar group variable. There is a useful
extension of the sign-code method, with additional explicit hypotheses.
It is a sufficient certificate, not a classification of arbitrary typed
kernels.

Let $\mathcal A,\mathcal B$ be finite type sets with probability measures
$\mu,\nu$. For each type pair take a real even function $w_{ab}:G\to\R$
whose Fourier coefficients $\lambda_{ab}(\xi)$ have the same positive
zero mode $\lambda_{ab}(0)=p$. The bipartite kernel is
\begin{equation}
 U((a,x),(b,y))=w_{ab}(x-y)
 \label{eq:typed-kernel}
\end{equation}
on $(\mathcal A\times G)\times(\mathcal B\times G)$ with product
probability measures. Let $S=-S$ be the union of its active nonzero
spectra, and put $D_L=\Delta_L(H)$, $D_R=\Delta_R(H)$.

\begin{theorem}[Typed sign-gauge certificate]
\label{thm:typed}
Suppose there are even sign functions
$s_a,t_b:S\to\{\pm1\}$ such that
\begin{equation}
 \sgn\lambda_{ab}(\xi)=s_a(\xi)t_b(\xi)
 \quad\text{whenever }\lambda_{ab}(\xi)\ne0,\quad\xi\in S.
 \label{eq:type-factor}
\end{equation}
Suppose every $s_a$ is $D_L$-safe and every $t_b$ is $D_R$-safe. Then
\begin{equation}
 t_{L,R}(H,U)\geq p^{e(H)}.
 \label{eq:typed-bound}
\end{equation}
Existence of such a certificate is equivalent to feasibility of a finite
affine linear system over $\F_2$.
\end{theorem}
\begin{proof}
Condition on all sampled vertex types. Expanding only the group
coordinates gives the conservation equations of Lemma~\ref{lem:flow},
with coefficient $\lambda_{a(u)b(v)}(\phi(uv))$ on each edge. For a
nonzero-weight flow, factor its sign using \eqref{eq:type-factor}.
The product of all $s$ factors groups into one zero-sum tuple at each
left vertex; each such tuple is positive by $D_L$-safety. The $t$
factors similarly group into positive tuples at the right vertices.
The full coefficient is therefore positive. For every type assignment
the zero flow has weight $p^e$, independent of the types. Averaging the
type assignments proves \eqref{eq:typed-bound}.

Choose inversion-orbit representatives $\xi_1,\ldots,\xi_m$ of $S$.
Let $z_a,u_b\in\F_2^m$ encode the proposed sign functions, and let
$b_{ab,j}$ be the known sign bit of a nonzero coefficient. Precisely the
required conditions are
\begin{equation}
 \begin{cases}
   z_{a,j}+u_{b,j}=b_{ab,j}
         &\text{for every active type--frequency pair},\\
   z_a\in\Rel_{D_L}(S)^\perp &\text{for every }a,\\
   u_b\in\Rel_{D_R}(S)^\perp &\text{for every }b.
 \end{cases}
 \label{eq:typed-linear-system}
\end{equation}
Theorem~\ref{thm:code} identifies the safety conditions with these
linear constraints. Thus \eqref{eq:typed-linear-system} is necessary and
sufficient for the stated certificate, and Gaussian elimination decides
its feasibility once relation-code bases have been supplied.
\end{proof}

At a single fixed frequency, \eqref{eq:type-factor} is possible if and
only if every cycle of the active bipartite type graph has positive sign
product. Necessity follows by multiplying around a cycle. For
sufficiency, choose a sign at one root of each component, propagate
along a spanning tree, and use the cycle condition to check every
remaining edge. The local additive constraints in
\eqref{eq:typed-linear-system} are extra compatibility requirements
across different frequencies.

The scalar theorem is recovered by using one type on each side,
assigning the spectral sign to one side and the constant positive sign
to the other. For general typed kernels, a negative total density
relative to $p^e$ implies that this sufficient certificate cannot exist;
it does not identify which type-cycle or additive constraint fails.
No claim is made that the rank-layer activation kernel in
\cite{OpenAIKernel} meets the real-even or sign-gauge hypotheses. Testing
that actual kernel requires additional work.

\section{Further research questions}
\label{sec:questions}

The following are precise questions left unresolved by this article.
They are proposed research targets, not assertions that each is new to
the entire literature.

\begin{enumerate}[leftmargin=*,label=\textbf{Q\arabic*.}]
\item \textbf{Scalar witnesses for the incidence pattern.}
Does there exist a finite abelian group and an even nonnegative scalar
kernel with $t(\Hstar,W)<p^{66}$? Theorem~\ref{thm:incidence} supplies a
necessary weighted negative-triple certificate. An affirmative answer
would need to show that negative contributions survive all positive
cycle and higher-flow contributions. A negative answer would separate
this scalar host class from more general typed or nonabelian hosts.

\item \textbf{The gap between sign safety and total positivity.}
For a fixed $D$, characterize the spectra for which Sidorenko holds for
all patterns of one-sided degree at most $D$, without requiring
termwise positivity. The theta example proves that the first failed
relation alone is insufficient. A useful intermediate target is a
criterion combining $D_3$, $Q_4$, and a finite collection of higher
moments for the incidence pattern.

\item \textbf{Sharper robust errors.}
Can one replace the relaxed one-sided error $\Err_H$ by an efficiently
computable functional retaining selected constraints on both sides?
For $\Hstar$, what is the best bound for negative full-flow mass in
terms of $D_3$, $M_3$, and low spectral moments? In particular, is a
substantially smaller replacement for the factor $22M_3^{21}$ available
under a natural boundedness or sparsity hypothesis?

\item \textbf{Weighted stability of the relation code.}
Under what spectral distribution hypotheses does small normalized
$D_D$ imply closeness, in a specified Fourier-weighted metric, to a
$D$-safe sign assignment? This is a weighted constraint-repair question.
The odd-characteristic and Boolean-circuit examples rule out a
hypothesis-free passage from bounded-degree safety to a global
translation character.

\item \textbf{Complete typed certificates.}
Can the affine system \eqref{eq:typed-linear-system} be weakened while
still certifying all full-flow weights, or can its failure be converted
into a bounded-size typed graph obstruction? A careful spectral analysis
of the source activation kernel is a natural test case. The distinction
between type-cycle frustration and frequency-relation frustration may
separate two different mechanisms for negative density.

\item \textbf{Efficient code filtrations in compressed groups.}
The shortest-witness algorithm is linear in the explicit cover size,
but $|K|$ may be exponential in a succinct group description. Develop
parameterized algorithms in spectrum size, rank, and degree that output
bases of $\Rel_D(S)$, preferably together with short integer-relation
certificates. The terminal code can also be studied through a relation
matrix and integer normal forms; extracting the complete degree
filtration is a finer problem.

\item \textbf{Other host classes and sharp constants.}
Determine which sharp constants survive restriction to simple $0$--$1$
Cayley graphs or extension to compact abelian groups. In the latter
case, absolute convergence or an approximation theorem is needed before
interchanging infinite Fourier sums. A nonabelian extension would need
to control matrix-valued Fourier products rather than merely signs.

\item \textbf{Formal verification and certified enumeration.}
Formalize the flow identity, the sign-code equivalence, the relaxed
negative-mass identity, and the injection from cycles to flows. For the
concrete pattern, kernel-check the pair and clique certificates and the
one-character polynomial. The integration plan distributed with this
article separates finite algebraic goals from analytic norm lemmas and
from computations that currently use a Python trusted boundary.
\end{enumerate}

\section{Reproducibility and proof status}
\label{sec:validation}

The package contains \texttt{article.tex}, its compiled PDF,
\texttt{verify.py}, the actual JSON output of the checker, machine-readable
face data, source and integration notes, and a build script. No downloaded
third-party PDF, font file, or entire repository is redistributed.

The checker uses only Python's standard library. Every asserted equality
or inequality uses integers or exact rational numbers; no floating-point
tolerance is used. It checks 56 graph--host instances across seven small
bipartite graphs and two Boolean groups, comparing a primal vertex sum
with a conserved-flow expansion. In those instances it also checks local
moment identities, the robust inventory bound, the girth surplus,
negative-mass control, the forest identity, and the cut-norm factor.

Additional tests compare tuple and integer-relation descriptions in 11
code cases, recover shortest witnesses in six cover-graph cases, and
check three exact theta polynomials with a negative term but no total
deficit. The incidence certificate verifies its degrees, pair
multiplicities, point-graph triangles, connectedness, short-cycle counts,
and the polynomial in Appendix~\ref{app:polynomial}.

The finite tests are not exhaustive over groups, functions, or graphs.
The general assertions are supported by the written proofs, not by an
inference from those tests. The typed extension has a written proof but
no dedicated typed-instance regression test in this package. No Lean
source is supplied as a purported proof, and no external peer or formal
verification is claimed.

\appendix
\section{The complete one-character polynomial of the incidence graph}
\label{app:polynomial}

Let $w_t$ be the two-point host \eqref{eq:rank-one}, and put $z=(t/p)^2$.
The density of $\Hstar$ is
\begin{equation}
 t(\Hstar,W_t)=p^{66}P_*(z),\qquad
 P_*(z)=\sum_{j=0}^{22}A_jz^j.
 \label{eq:Pstar}
\end{equation}
Here $A_j$ counts even-degree edge subsets of $\Hstar$ with $2j$ edges.
The degree cannot exceed 22: a face vertex has degree three, so an
even-degree selected subgraph uses at most two of its incidences,
for at most 44 selected edges in total.

There is a small independent vertex-side formula. For a sign assignment
$s=(s_0,\ldots,s_{12})\in\{\pm1\}^{13}$ let $k(s)$ be the number of
monochromatic faces and put
\begin{equation}
 N_k=\#\{s\in\{\pm1\}^{13}:k(s)=k\}.
 \label{eq:Nk}
\end{equation}
Averaging the spin at a face $\{a,b,c\}$ gives
\[
 \E_{\tau=\pm1}(p+ts_a\tau)(p+ts_b\tau)(p+ts_c\tau)
   =p^3\bigl(1+z(s_as_b+s_as_c+s_bs_c)\bigr).
\]
The parenthesized pair sum is $3$ for a monochromatic face and $-1$
otherwise. Hence the exact identity
\begin{equation}
 P_*(z)=2^{-13}\sum_{k=0}^{22}N_k(1+3z)^k(1-z)^{22-k}
 \label{eq:Pstar-primal}
\end{equation}
computes the entire polynomial by enumerating only $2^{13}=8192$
point-spin assignments. Equivalently,
\begin{equation}
 A_j=2^{-13}\sum_{k=0}^{22}N_k
       \sum_{a+b=j}\binom ka3^a\binom{22-k}{b}(-1)^b,
 \label{eq:Aj}
\end{equation}
where binomial coefficients outside their natural range vanish.
The values obtained by exact enumeration are in Table~\ref{tab:poly}.
They begin
\[
 P_*(z)=1+33z^2+88z^3+823z^4+4232z^5+\cdots.
\]
The zero linear coefficient and the coefficients $33,88$ agree with
the independent girth and short-cycle counts. Also
\begin{equation}
 P_*(1)=\sum_j A_j=2^{32},
 \label{eq:Pstar-one}
\end{equation}
in agreement with the cycle-space dimension
$66-35+1=32$. To see the dimension formula directly, choose a spanning
tree. Any choice of the 32 chord bits determines a unique even-degree
labeling on the tree by successive leaf elimination, including the
last root constraint. Thus there are exactly $2^{32}$ such subgraphs.

The coefficient table is an exact finite certificate accompanied by an
executable derivation, not a formally kernel-checked enumeration.
\begin{table}[p]
\centering
\begin{tabular}{rrr}
\toprule
$j$ & $N_j$: point-spin assignments & $A_j$: even subgraphs with $2j$ edges\\
\midrule
0 & $96$ & $1$ \\
1 & $352$ & $0$ \\
2 & $696$ & $33$ \\
3 & $1\,012$ & $88$ \\
4 & $1\,130$ & $823$ \\
5 & $1\,176$ & $4\,232$ \\
6 & $1\,080$ & $24\,831$ \\
7 & $784$ & $127\,584$ \\
8 & $646$ & $603\,274$ \\
9 & $492$ & $2\,540\,960$ \\
10 & $290$ & $9\,417\,578$ \\
11 & $180$ & $30\,237\,520$ \\
12 & $102$ & $83\,066\,126$ \\
13 & $80$ & $192\,776\,048$ \\
14 & $38$ & $373\,654\,910$ \\
15 & $8$ & $599\,077\,600$ \\
16 & $10$ & $785\,645\,157$ \\
17 & $12$ & $830\,015\,520$ \\
18 & $6$ & $690\,493\,733$ \\
19 & $0$ & $436\,109\,080$ \\
20 & $0$ & $196\,772\,043$ \\
21 & $0$ & $56\,595\,016$ \\
22 & $2$ & $7\,805\,139$ \\
\bottomrule
\end{tabular}
\caption{Exact enumeration for the incidence pattern. The second column
sums to $8192$, and the third to $2^{32}$. The $N_j$ column, together
with Equation~\eqref{eq:Aj}, reproduces every entry of the $A_j$ column.}
\label{tab:poly}
\end{table}

\clearpage

\section{A finite-group normalization pitfall}
\label{app:normalization}

A general integer matrix that has full row rank over $\Q$ need not
define a surjective map on an arbitrary finite abelian group. Thus
Fourier formulas normalized by a putative $|G|^{k-m}$ kernel size need
a surjectivity hypothesis or an explicit multiplicity factor.

For a small example, let $G=\Z/2\Z$ and
\[
 L=\begin{pmatrix}1&1\\1&-1\end{pmatrix}.
\]
This matrix has determinant $-2$ over $\Z$ but its two rows coincide on
$G$. Its kernel has size two, not $|G|^{2-2}=1$. For $f\equiv1$ a
uniform kernel average is one, whereas summing the corresponding
Fourier coefficient products over both row frequencies gives two:
the map on row frequencies has a two-element kernel.

This does not invalidate the special graphic identities needed here.
It explains why Lemma~\ref{lem:flow} derives the expansion by averaging
\emph{vertex variables} directly. The theorem then has no hidden
matrix-rank assumption and works equally in even and odd characteristic.
This normalization issue is separate from the stronger graphical
structure used in subdivision arguments such as \cite{YuqiZhao}.

\section{Source snapshot and integration boundary}
\label{app:sources}

The repository references below were inspected on 7 October 2026. The
OpenAI repository head was
\begin{center}
\small\texttt{adc7f1241b42e322a6451854ab7e4b4c146bf78a},
\end{center}
and the ProveIt head was
\begin{center}
\small\texttt{a41aa448d67395387593f1931cda142836afccdc}.
\end{center}
The finite-complex source and ProveIt definitions were also read at
these explicit commits. Other selected material was read from the moving
branch, with source blob identifiers recorded where returned by the
connector. The review was targeted, not an audit of all manuscripts or
all formalization files in either repository.

For integration, use a separate research directory, for example
\nolinkurl{Combinatorics/Sidorenko/Research/LocalFourierSignCodes/}, adapting
that path to the maintainer's chosen taxonomy. The mathematical APIs
should begin with normalized finite-group Fourier coefficients and
explicit bipartitions. Existing unnormalized Gowers definitions require
the conversion stated after \eqref{eq:Q4-U2}. No statement catalogue
should be imported as an axiom establishing a theorem.

The provided \texttt{INTEGRATION.md} lists the proposed formalization
order and invariant checks. The \texttt{SOURCE\_AUDIT.md} records the
limited source review, priority qualifications, and typed-versus-scalar
boundary. There is no claim that the repository has been modified.

\begin{thebibliography}{9}
\bibitem{OpenAICatalogue}
OpenAI, \emph{Mathematics manuscript collection},
\texttt{openai/math}, \texttt{CONTENTS.md}, especially result family 161.
Repository snapshot consulted 7 October 2026.
\url{https://github.com/openai/math}.

\bibitem{OpenAIComplex}
OpenAI mathematics manuscript collection,
\emph{A counterexample to Sidorenko's conjecture}, dated
23 September 2026, finite-complex source section.
Path under \texttt{openai/math}:\par
\nolinkurl{preprints/A-counterexample-to-Sidorenkos-conjecture-September-23-2026/build/sections/complex.tex}.
Source blob \texttt{178b0e9da21867a727ea847dae8e537b9d719a70}.
The explicit face data, rather than the claimed counterexample theorem,
are the input used here.

\bibitem{OpenAIKernel}
Same manuscript, kernel-construction source section,\par
\nolinkurl{preprints/A-counterexample-to-Sidorenkos-conjecture-September-23-2026/build/sections/kernel.tex}.
Source blob \texttt{bbe11be8fffe968012060ec43e6aff597e922831}.
The inspected portion defines the rank-layer and typed activation
kernels; the complete analytic proof was not audited here.

\bibitem{ProveItDefinitions}
Vladimir Reshetnikov, \emph{ProveIt},
Gowers--Szemeredi finite-combinatorial definitions,\par
\nolinkurl{Combinatorics/Ramsey/Lean/GowersSzemeredi/Definitions.lean}.
Pinned read at the ProveIt commit recorded in Appendix~\ref{app:sources}.
\url{https://github.com/VladimirReshetnikov/ProveIt}.

\bibitem{ProveItAudit}
\emph{ProveIt}, targeted fixed-radius Bohr source audit,\par
\nolinkurl{Combinatorics/Ramsey/Research/GowersSzemeredi/local-quantitative-refinements/30-fixed-radius-bohr-SOURCE_AUDIT.md}.
Source blob \texttt{ea818e377eaf18b512c40e0081674ed78e5197ad}.
Consulted for normalization and verification-status context, not as a
premise in the present proofs.

\bibitem{YuqiZhao}
Yuqi Zhao, \emph{Sidorenko-Type Inequalities for Even Subdivisions over
Finite Abelian Groups}, arXiv:2507.15723v1, 21 July 2025.
\url{https://arxiv.org/abs/2507.15723v1}.

\bibitem{PositiveGraphs}
David Conlon, Joonkyung Lee, and Leo Versteegen,
\emph{Around the positive graph conjecture}, arXiv:2404.17467, 2024.
\url{https://arxiv.org/abs/2404.17467}.
The abstract was consulted for the distinction between arbitrary-host
positivity and the restricted Fourier certificates considered here.
\end{thebibliography}

\end{document}
